\section{The two lines}\label{sec:twolines}

The author proposed a second line at $\re s=-\tfrac12$ that mirrors the critical line, and asked for it to be built by shifting the whole line, as in the programme's Hurwitz shift, and for the anomaly, chirality and holonomy language of that picture to be made exact.
This section gives the exact objects. (The Hurwitz flow does not translate the zero set; what vanishes on the shifted set is its first jet, Theorem~\ref{thm:jet}.) The typed matching of the author's terms with the objects is a proposal (Section~\ref{sec:bridges}); the objects and the facts below are theorems, none of them assumes RH, and Proposition~\ref{prop:races} states its GRH hypotheses where it uses them.
Recall $A(s)=-\bar s$ (the mirror through the pole at $0$), $s^\#=1-\bar s$ (the reflection through $\tfrac12$) and $T_c(s)=s+c$; a direct computation gives $s^\#=T_1(A(s))$.

\subsection{The first Hurwitz jet and the mirror set}

\begin{theorem}[the first Hurwitz jet]\label{thm:jet}
\begin{enumerate}[label=(\alph*)]
\item For $s\ne1$, $\partial_t\zeta(s,1+t)|_{t=0}=-s\zeta(s+1)$, where at $s=0$ the right side means its holomorphic extension; more generally $\partial_t^m\zeta(s,1+t)|_{t=0}=(-1)^m(s)_m\zeta(s+m)$.
\item $s\zeta(s+1)$ is entire and equals $1$ at $s=0$; so the first jet equals $-1$ at $s=0$.
\item In $-1<\re s<0$ the zeros of the first jet are exactly $\Zs-1$, with multiplicities; in $-m<\re s<1-m$ the zeros of the $m$-th jet are exactly $\Zs-m$.
\item $\zeta$ has no zero on $\Zs-1$.
\end{enumerate}
\end{theorem}
\status{Proved here (claude-ab, \note{17\_} Theorem 17.1; part (a) on $\re s>1$ is the programme's shift-flow theorem). Checked numerically (the $t$-derivative at $\rho_k-1$ is below $4\cdot10^{-29}$ for $k=1,2,3$). Refereed (fifth pass).}

\begin{proof}
(a) On $\re s>1$, $\partial_a\sum_n(n+a)^{-s}=-s\sum_n(n+a)^{-s-1}$ termwise. The Euler--Maclaurin representation shows that $\zeta(s,a)$ is jointly analytic on $\{s\ne1\}\times\{\re a>0\}$, so both sides are analytic in $s\ne1$ and agree everywhere by the identity theorem; the higher jets follow in the same way.
(b) $\zeta(s+1)$ has a simple pole at $s=0$ with residue $1$.
(c) On $-1<\re s<0$ the factor $s$ does not vanish and $0<\re(s+1)<1$, so the zeros are the $s$ with $s+1\in\Zs$; the trivial zeros lie outside. For $m\ge2$, $(s)_m$ vanishes only at $0,-1,\dots,1-m$, none of real part in $(-m,1-m)$.
(d) Points of $\Zs-1$ have real part in $(-1,0)$, where $\zeta$ has no zeros ($\zeta\ne0$ on $\re s\ge1$, and the functional equation leaves only the trivial zeros in $\re s\le0$).
\end{proof}

Implicit differentiation of $\zeta(\rho(t),1+t)=0$ gives the velocity $\rho'(0)=\rho\zeta(\rho+1)/\zeta'(\rho)$ of a simple zero, read on the line $\re s=1+\re\rho$; for $\rho_1$ it is $0.98358+9.65780i$, so the first zero leaves the critical line to the right as $t$ increases.

\begin{theorem}[the shifted set is the mirror set]\label{thm:mirror}
$A(\Zs)=\Zs-1$ as sets with multiplicities, unconditionally. For every $s$, $A(s)-(s-1)=1-2\re s$; hence $A(\rho)=\rho-1$ exactly when $\re\rho=\tfrac12$, and
\[
\text{RH}\iff A(\rho)=\rho-1\ \text{for every nontrivial zero }\rho.
\]
Moreover $A(\rho)=\rho^\#-1$: the mirror of $\rho$ is the shift of its partner.
\end{theorem}
\status{Proved here (claude-ab, \note{17\_} Theorem 17.2). Refereed (fifth pass).}

\begin{proof}
$-\bar\rho=(1-\bar\rho)-1$, and $1-\bar\rho\in\Zs$ with the multiplicity of $\rho$; $-\bar s-(s-1)=1-(s+\bar s)$.
\end{proof}

\begin{figure}[t]
\centering
\includegraphics[width=\linewidth]{fig_mirror_line_double_helix.pdf}
\caption{Left: the critical line, its mirror $A(\Zs)=\Zs-1$ on $\re s=-\tfrac12$ (the nodal set of the first Hurwitz jet in $-1<\re s<0$), and a \emph{hypothetical} off-line zero, drawn only to show that its mirror rung and its shift rung differ; no such zero exists below height $3\cdot10^{12}$. Right: on the Clifford torus the two lines form one $(1,-1)$ circle $\Gamma$, which meets the $(1,1)$ fibre $D$ only at $t=\pm\infty$ (\note{17\_} Proposition 17.3).}
\label{fig:mirror}
\end{figure}

So the author's two operations produce the same line: as sets the mirror and the shift agree unconditionally, as maps on the zeros they agree exactly under RH, and the set carries no zero of $\zeta$.

\begin{proposition}[the blind lemma]\label{prop:blind}
Let $H_\sigma=L^2((0,\infty),u^{2\sigma-1}du)$ with $U_af(u)=f(au)$. If $\sigma\ne\sigma'$, every bounded $K:H_\sigma\to H_{\sigma'}$ with $KU_a=U_aK$ for all $a>0$ is zero. The inversion $f\mapsto f(1/u)$ is an isometry of $H_\sigma$ onto $H_{-\sigma}$ acting on Mellin transforms by $s\mapsto-s$ and reversing the flow.
\end{proposition}
\status{Proved here (claude-ab, \note{17\_} Proposition 17.4). Refereed (fifth pass).}

\begin{proof}
$\|U_af\|_\sigma=a^{-\sigma}\|f\|_\sigma$, so $U_a=a^{-\sigma}V_a$ with $V_a$ unitary on $H_\sigma$, and likewise on $H_{\sigma'}$. Then $V'_aK=a^{\sigma'-\sigma}KV_a$, so $\|Kf\|\le a^{\sigma'-\sigma}\|K\|\|f\|$ for all $a>0$; let $a\to\infty$ or $a\to0$. The inversion statement is the substitution $v=1/u$.
\end{proof}

Boundedness is essential: the identity on $C_c^\infty(0,\infty)$ is a nonzero closable densely defined operator $H_{1/2}\to H_{-1/2}$ commuting with every $U_a$.

\subsection{The two centres of the explicit formula}

With $(g_1\star g_2)(x)=\int g_1(y)g_2(x/y)\,dy/y$ and $\hat g(s)=\int_0^\infty g(x)x^{s-1}dx$, use Weil's explicit formula in Bombieri's form \cite{Bombieri}.

\begin{proposition}[two centres]\label{prop:centres}
(a) With $g^*(x)=x^{-1}\cl{g(1/x)}$ and $h=g\star g^*$: $\hat h(s)=\hat g(s)\cl{\hat g(1-\bar s)}$, both sides of the explicit formula are real, an on-line zero contributes $m_\rho|\hat g(\rho)|^2\ge0$, and an off-line pair $\{\rho,\rho^\#\}$ contributes $2m_\rho\re(\hat g(\rho)\cl{\hat g(\rho^\#)})$, an indefinite form with eigenvalues $\pm m_\rho$ on $(\hat g(\rho),\hat g(\rho^\#))$.
(b) With $g^\flat(x)=\cl{g(1/x)}$ and $h=g\star g^\flat$: $\hat h(s)=\hat g(s)\cl{\hat g(-\bar s)}$, the prime term has imaginary part $\sum_n\Lambda(n)(1-1/n)\im h(n)$, and the sum over zeros is not real in general.
\end{proposition}
\status{Proved here (claude-ab, \note{17\_} Proposition 17.5). Both sides computed independently for the test $g(x)=\exp(-(\log x)^2/2+\tfrac i2\log x-i\gamma_1\log x)$: centre $\tfrac12$ gives $2\pi$ on both sides to $4\cdot10^{-15}$; centre $0$ gives $2\pi e^{i/2}$ to $1.6\cdot10^{-12}$. Refereed (fifth pass).}

\begin{proof}
Substituting $y=1/x$ gives $\widehat{g^*}(s)=\cl{\hat g(1-\bar s)}$ and $\widehat{g^\flat}(s)=\cl{\hat g(-\bar s)}$; the Mellin transform turns $\star$ into a product. The zero side of (a) is real because $\rho\mapsto\rho^\#$ permutes $\Zs$ with multiplicities; $h(1/x)=x\cl{h(x)}$ makes $h(n)+h(1/n)/n=2\re h(n)$. In (b), $h(1/x)=\cl{h(x)}$, so $h(n)+h(1/n)/n=h(n)+\cl{h(n)}/n$, whose imaginary part is $(1-1/n)\im h(n)$.
\end{proof}

So a real partner pairing selects the centre $\tfrac12$ and excludes the pole at $0$; with that centre, the only possible source of negativity is the hyperbolic block of an off-line pair, and by Weil's criterion \cite{Weil,Bombieri} a negative value exists if and only if RH fails.

\subsection{The two-line product and its symmetries}

\begin{proposition}[the two-line product on the line of the pole at $0$]\label{prop:G}
$G(s)=\xi(s)\xi(s+1)/16$ satisfies $G(-\bar s)=\cl{G(s)}$ and $G(it)=|\xi(1+it)|^2/16>0$ for real $t$. Its zero set is $\Zs\sqcup(\Zs-1)$.
\end{proposition}
\status{Proved here (claude-ab, \note{22\_} Proposition 22.4, \note{24\_} Lemma 24.3; register S37). Checked to $10^{-30}$. Refereed (eighteenth pass).}

\begin{proof}
From $\xi(1-s)=\xi(s)$ and $\xi(\bar s)=\cl{\xi(s)}$: $\xi(-\bar s)=\xi(1+\bar s)=\cl{\xi(1+s)}$ and $\xi(1-\bar s)=\cl{\xi(s)}$, which gives the first identity; at $s=it$, $\xi(it)=\xi(1-it)=\cl{\xi(1+it)}$, and $\xi(1+it)\ne0$ because $\zeta$ has no zeros on $\re s=1$.
\end{proof}

The maps $\mathrm{id}$, $A$, $P$ ($=\#$ on $\Zs$ and $s\mapsto-1-\bar s$ on $\Zs-1$) and $T=AP$ ($s\mapsto s-1$ on $\Zs$, $s\mapsto s+1$ on $\Zs-1$) form a Klein four-group preserving $\Zs\sqcup(\Zs-1)$ with multiplicities. An on-line zero has orbit $\{\rho,\rho-1\}$; an off-line zero has the orbit of four points $\{\rho,\rho^\#,\rho^\#-1,\rho-1\}$ (the author's doublet). The loop ``across by $A$, back by $T^{-1}$'' has holonomy $T^{-1}A=\#$ on $\Zs$, trivial exactly at on-line zeros; and the zeros of $G$ in the open strip $-\tfrac12<\re s<\tfrac12$ (the bulk between the lines) are exactly the zeros with $\re\rho<\tfrac12$ and their mirrors, so the bulk is empty if and only if RH holds (\note{22\_} Propositions 22.6--22.7; the checks are elementary).

\subsection{The chiral cross form is Weil's form on the diagonal}

For finitely supported $u$ on $\Zs$ let $\Omega(u)=\sum_\rho m_\rho u(\rho)\cl{u(\rho^\#)}$ (Weil's pairing), and for $f$ on $\Zs\sqcup(\Zs-1)$ let $Q_A(f)=\sum_wm_wf(w)\cl{f(Aw)}$ (the chiral cross form, centred at $0$). Both are real because $\#$ and $A$ are involutions preserving multiplicities.

\begin{proposition}[the chiral form and Weil's form]\label{prop:chiral}
Let $f$ be finitely supported on $\Zs\sqcup(\Zs-1)$, or the restriction of an element of $\B$ (then all sums converge absolutely). Put $u(\rho)=\tfrac12(f(\rho)+f(\rho-1))$ and $v(\rho)=\tfrac12(f(\rho)-f(\rho-1))$, the $T$-even and $T$-odd parts. Then
\[
Q_A(f)=2\Omega(u)-2\Omega(v).
\]
\begin{enumerate}[label=(\alph*)]
\item $A$ has no fixed point on $\Zs\sqcup(\Zs-1)$, so $Q_A$ is an orthogonal sum of hyperbolic planes, of signature $(n,n)$ on data supported on $n$ pairs $\{w,Aw\}$, whether or not RH holds.
\item $\Omega(u)=\sum_{\text{on-line}}m_\rho|u(\rho)|^2+\sum_{\text{off-line pairs}}2m_\rho\re(u(\rho)\cl{u(\rho^\#)})$; so $\Omega\ge0$ on all finitely supported $u$ exactly when every zero is on the line.
\item Under RH the $T$-even data form a maximal positive subspace for $Q_A$; if RH fails, they contain a negative direction.
\end{enumerate}
\end{proposition}
\status{Proved here (claude-ab, \note{42\_} Proposition 42.2; register S72). Checks B1--B5. Refereed (nineteenth pass, which extended it to data from $\B$).}

\begin{proof}
For $\rho\in\Zs$, $A\rho=\rho^\#-1\in\Zs-1$ and $A(\rho-1)=\rho^\#$, so $Q_A(f)=\sum_\rho m_\rho[f(\rho)\cl{f(\rho^\#-1)}+f(\rho-1)\cl{f(\rho^\#)}]$. Substituting $f=u+v$ on $\Zs$ and $f=u-v$ on $\Zs-1$, the cross terms cancel in pairs and $2\sum_\rho m_\rho[u(\rho)\cl{u(\rho^\#)}-v(\rho)\cl{v(\rho^\#)}]$ remains.
(a) $Aw=w$ means $\re w=0$, where $G$ has no zero (Proposition~\ref{prop:G}); on a pair $\{w,Aw\}$ the form is $2m_w\re(f(w)\cl{f(Aw)})$.
(b) $\#$ fixes the on-line zeros and pairs the off-line ones; for an off-line $\rho$ take $u=1$ at $\rho$ and $-1$ at $\rho^\#$.
(c) Under RH, $\Omega(u)=\sum m|u|^2$; a subspace strictly containing the $T$-even data contains a nonzero $T$-odd vector $x$, and $Q_A(x)=-2\Omega(v)<0$.
\end{proof}

\begin{corollary}[an RH criterion in two-line terms]\label{cor:chiral}
For $k\in\B$ put $f_k(s)=F_*(s+1)F_*(s-1)k(s)+F_*(s+2)F_*(s)k(s+1)$. Then $f_k\in\B$ is $T$-even on the zero set, with $u(\rho)=F_*(\rho+1)F_*(\rho-1)k(\rho)$, and
\[
\text{RH}\iff Q_A(f_k)\ge0\ \text{for every }k\in\B .
\]
\end{corollary}
\status{Proved here (claude-ab, \note{42\_} Corollary 42.3). Equivalent to RH by construction: its test data can be localised at the zeros. Refereed (nineteenth pass).}

\begin{proof}
$F_*(\rho)=0$ gives $f_k(\rho)=F_*(\rho+1)F_*(\rho-1)k(\rho)=f_k(\rho-1)$. $F_*(\rho\pm1)\ne0$, since $\re(\rho\pm1)\notin(0,1)$ and $F_*$ vanishes only on $\Zs$. Finite sums of the programme's section representatives realise any finitely supported values of $k$ on $\Zs$, so $u$ ranges over all finitely supported data and Proposition~\ref{prop:chiral}(b) applies. Under RH, $\Omega(u)=\sum m|u|^2\ge0$.
\end{proof}

$Q_A(f)$ is Weil's explicit-formula functional of $G$ applied to $h=f\cdot f^A$, $f^A(s)=\cl{f(-\bar s)}$: for $a>1$, $Q_A(f)=(2\pi i)^{-1}\int_{(a)}[h(s)+h(-s)](G'/G)(s)\,ds$, where on $\re s=a$
\[
\frac{G'}{G}(s)=\frac2s+\frac1{s-1}+\frac1{s+1}-\log\pi+\tfrac12\psi(\tfrac s2)+\tfrac12\psi(\tfrac{s+1}2)-\sum_n\Lambda(n)(1+n^{-1})n^{-s}
\]
(the horizontal sides vanish along the good heights of Lemma~\ref{lem:goodheights}; checked to 25 digits). Whether the diagonal family $\{f_k\}$ is easier to control on the prime side than Weil's family is open.

\subsection{Counting: parity and the holonomy index}

\begin{proposition}[parity sees only the count on the line]\label{prop:parity}
For $T>0$ not an ordinate, $(-1)^{N(T)}=(-1)^{N_0(T)}=-\sgn Z(T)$, with no simplicity assumption. In the two-line system the count up to height $T$ is $n_G(T)=2N_0(T)+4P(T)$; its reduction mod $2$ vanishes identically, and $i^{n_G(T)}=(-1)^{N_0(T)}=-\sgn Z(T)$. Neither depends on $P(T)$; mod $8$, $n_G$ gives $P(T)\bmod2$ only once $N_0(T)\bmod4$ is known, and $(N_0,P)=(0,1)$ and $(2,0)$ both give $n_G=4$.
\end{proposition}
\status{Proved here (claude-ab, \note{22\_} 22.9, corrected in the eighteenth referee pass and in the final pass). Checked at eight heights up to $1000$.}

\begin{proof}
With $s=\tfrac12+it$, $s(s-1)=-(\tfrac14+t^2)$ and $\pi^{-s/2}\Gamma(s/2)\zeta(s)=\pi^{-1/4}|\Gamma(\tfrac14+\tfrac{it}2)|Z(t)$, so $\xi(\tfrac12+it)=-\tfrac12(\tfrac14+t^2)\pi^{-1/4}|\Gamma(\tfrac14+\tfrac{it}2)|\,Z(t)$ is real and $\sgn\xi(\tfrac12+iT)=-\sgn Z(T)$.
$\xi(\tfrac12)>0$, and $\xi(\tfrac12+it)$ changes sign exactly at on-line zeros of odd order, whose number has the parity of $N_0(T)$; off-line zeros come in pairs of equal height, so $N(T)\equiv N_0(T)\pmod2$.
The zeros of $G$ up to height $T$ are those of $\xi$ and their translates, so $n_G=2N(T)=2N_0+4P$.
\end{proof}

\begin{proposition}[the holonomy index is Turing's]\label{prop:index}
Let $P(T)$ be the number, with multiplicity, of zeros with $0<\gamma\le T$ and $\re\rho<\tfrac12$ (one member of each off-line pair). Then
\[
I(T):=\frac{N(T)-V(T)}2=P(T)+\sum_{\substack{\re\rho=1/2\\0<\gamma\le T}}\lfloor m_\rho/2\rfloor .
\]
So $I(T)\ge0$, and $I(T)=0$ exactly when every zero up to height $T$ is on the line and simple; a positive value does not separate off-line pairs from multiple zeros on the line.
\end{proposition}
\status{Proved here (claude-ab, \note{42\_} Proposition 42.5). Checks D1--D4 and R5 (for $\zeta$: $N=V=29$ at $T=100$, $N=V=108$ at $T=250$). Refereed (nineteenth pass).}

\begin{proof}
Off-line pairs contribute $2P(T)$ to $N(T)$. $Z$ is real-analytic and real on $\R$, its zeros are the on-line zeros with the same orders, and it changes sign exactly at those of odd order, so $V(T)=\sum_{\text{on-line}}(m_\rho-2\lfloor m_\rho/2\rfloor)$. Subtract.
\end{proof}

$I(T)$ is computed from boundary data ($N(T)=\theta(T)/\pi+1+S(T)$ by the argument principle, $V(T)$ from the line), and $I(T)=0$ is exactly what verifications of RH by Turing's method \cite{Turing} establish up to a height; RH has been verified up to $3\cdot10^{12}$ (Platt--Trudgian \cite{PlattTrudgian}). The statement $I(T)=0$ for all $T$ is RH together with the simplicity of all zeros.

\begin{proposition}[the second line and prime races]\label{prop:races}
For a primitive character $\chi$ mod $q$ and $G_\chi(s)=\Lambda(s,\chi)\Lambda(s+1,\chi)$, the prime coefficients of $-G_\chi'/G_\chi$ are $\Lambda(n)\chi(n)(1+n^{-1})$. For nonprincipal $\chi$, $\sum_{p\le x}\chi(p)\log p/p$ converges; hence the two-line races mod $4$ differ from the one-line races by convergent series, which vanish after Rubinstein--Sarnak's normalisation, and the two-line races have the same limiting logarithmic distribution (under GRH) and the same logarithmic density $\delta(4;3,1)\approx0.9959$ (under GRH and linear independence) as the one-line races \cite{RubinsteinSarnak,Devin}.
\end{proposition}
\status{Proved here (claude-ab, \note{42\_} Proposition 42.4; partial summation with the prime number theorem for progressions, Davenport \cite{Davenport}, Ch.~20). Checks C1--C3 up to $10^7$. Refereed (nineteenth pass).}

\subsection{The separator}

Let $I_-$ be the ideal of functions vanishing to full order on $\Zs-1$.

\begin{lemma}[the separator lemma with matched rates]\label{lem:separator}
Let $G\in\B$ have zero divisor $D_+\sqcup D_-$ with $D_+\subset\{\re s>0\}$ and $D_-\subset\{\re s<0\}$. Suppose there are $\alpha\ge0$, $\varepsilon>0$ and $N,K\ge0$ such that, with $r(t)=\varepsilon(1+t^2)^{-N}$,
(a) $|G(x+it)|\le C_A(1+|t|)^{k_A}e^{-\alpha|t|}$ on every strip $|x|\le A$, and
(b) $G$ has no zeros in the corridor $\mathcal{K}=\{|x|\le r(t)\}$ and $|G|\ge c(1+|t|)^{-K}e^{-\alpha|t|}$ there.
Then there is an entire $E$ of polynomial growth on every vertical strip with $E-1$ vanishing to full order on $D_+$ and $E$ vanishing to full order on $D_-$; consequently $\B/(I_{D_+}\cap I_{D_-})\to\B/I_{D_+}\oplus\B/I_{D_-}$ is a topological isomorphism with inverse $([f],[g])\mapsto[Ef+(1-E)g]$, commuting with multiplication by $s$ and by every $a^s$.
\end{lemma}
\status{Proved in the programme for $\zeta$ (\lbl{GSL4--GSL5}); stated with general constants and checked by claude-ab (\note{24\_} Lemma 24.1). Standard $\bar\partial$ technique. Refereed (eighth pass).}

\emph{Outline.} Put $\chi(x+it)=\eta(x/r(t))$ with a cutoff $\eta$ switching from $0$ to $1$; then $|\bar\partial\chi|\le C(1+|t|)^{2N}$, and $h=\bar\partial\chi/G$ is supported in $\mathcal{K}$ with $|h|\le C(1+|t|)^{2N+K}e^{\alpha|t|}$ by (b). The solution $u=\pi^{-1}\int e^{(s-w)^2}(s-w)^{-1}h(w)\,dA(w)$ of $\bar\partial u=h$ satisfies $|u|\le C_Ae^{\alpha|t|}(1+|t|)^{2N+K}$, because the Gaussian dominates $e^{\alpha|t-v|}$ and the singularity $1/|s-w|$ is integrable over the thin corridor. So $E=\chi-Gu$ is entire, and by (a) the two exponentials cancel exactly: $|Gu|$ grows polynomially. Near $D_\pm$, $h=0$ and $E-\chi=-Gu$ vanishes to full order. If the lower rate in (b) were $e^{-\alpha_2|t|}$ with $\alpha_2>\alpha$, the method would give only $|Gu|\le C(1+|t|)^ke^{(\alpha_2-\alpha)|t|}$. \qed

For $\zeta$, $G=F_*(s)F_*(s+1)=s^2(s^2-1)\zeta(1-s)\zeta(1+s)/(64\cos(\pi s/2))$ and $\alpha=\pi/2$: the only exponential factor, $1/\cos(\pi s/2)$, has the same size on every vertical line, which is why the rates match.

\begin{proposition}[the separator for every primitive Dirichlet $L$-function]\label{prop:separatorL}
Let $\chi$ be primitive mod $q\ge3$ with $\chi(-1)=(-1)^a$, $\Lambda_\chi(s)=(q/\pi)^{(s+a)/2}\Gamma(\tfrac{s+a}2)L(s,\chi)$, $Z_\chi$ its zero divisor and $I_\chi$, $I_{\chi,-}$ the ideals of full-order vanishing on $Z_\chi$ and $Z_\chi-1$. There is an entire $E_\chi$ of polynomial growth on strips that is $1$ to full order on $Z_\chi$ and $0$ to full order on $Z_\chi-1$, so $\B/(I_\chi\cap I_{\chi,-})\cong\B/I_\chi\oplus\B/I_{\chi,-}$ equivariantly.
\end{proposition}
\status{Proved here (claude-ab, \note{24\_} Proposition 24.4; register S44). Matched rates checked on $\re s=0$ for two characters. Refereed (eighth pass). Novelty: not searched.}

\begin{proof}
We verify Lemma~\ref{lem:separator} for $G_\chi(s)=\Lambda_\chi(s)\Lambda_\chi(s+1)$ with $\alpha=\pi/2$, using (F1) $\Lambda_\chi(s)=W(\chi)\Lambda_{\bar\chi}(1-s)$ with $|W(\chi)|=1$, (F2) $L(1+it,\chi)\ne0$, and (F3) Stirling's formula (Davenport \cite{Davenport}, Chs.~9, 4, 14, 10).
\emph{Growth.} $S(x)=\sum_{n\le x}\chi(n)$ is bounded and $q$-periodic; $L(s,\chi)=s\int_1^\infty S(x)x^{-s-1}dx$ on $\re s>0$, and repeated integration by parts against periodic antiderivatives of mean zero continues $L$ with $|L|+|L'|\le C_{A,q}(1+|t|)^K$ on $|\re s|\le A$; on $\tfrac34\le\sigma\le2$, $|L'(s,\chi)|\le C_1(1+|t|)$. By (F3), $\Lambda_\chi\in\B$.
\emph{Separation.} By the Euler product and (F2), $L\ne0$ on $\re s\ge1$; by (F1), $\Lambda_\chi\ne0$ on $\re s\le0$. So $Z_\chi\subset\{0<\re s<1\}$ and $Z_\chi-1\subset\{-1<\re s<0\}$.
\emph{Upper rate.} By (F1), $G_\chi(s)=W(\chi)(q/\pi)^{1+a}\Gamma(\tfrac{1+a-s}2)\Gamma(\tfrac{1+a+s}2)L(1-s,\bar\chi)L(1+s,\chi)$, and the Gamma product is $\pi/\cos(\pi s/2)$ for $a=0$ and $(\pi s/2)/\sin(\pi s/2)$ for $a=1$, of modulus at most $C_A(1+|t|)e^{-\pi|t|/2}$ on strips.
\emph{Lower rate.} For $\sigma>1$, $\log\big(\zeta(\sigma)^3|L(\sigma+it,\chi)|^4|L(\sigma+2it,\chi^2)|^2\big)=\sum_p\sum_kk^{-1}p^{-k\sigma}[3+4\re(\chi(p)^kp^{-ikt})+2\re(\chi(p)^{2k}p^{-2ikt})]\ge0$, each bracket being $(1+2\cos\theta)^2$ or $3$. With $\zeta(1+\delta)\le2/\delta$ and $|L(1+\delta+2it,\chi^2)|\le C(1+|t|)$ this gives $|L(1+\delta+it,\chi)|\ge c\delta^{3/4}(1+|t|)^{-1/2}$, and with the derivative bound, $|L(1+x+it,\chi)|\ge c(1+|t|)^{-5}$ for $|x|\le c(1+|t|)^{-6}$; the same holds for $L(1-x-it,\bar\chi)$, and the Gamma product is at least $ce^{-\pi|t|/2}$ there. For $|t|\le1$, $G_\chi(it)=\Lambda_\chi(it)\Lambda_\chi(1+it)\ne0$ by (F1)--(F2), and compactness gives a zero-free neighbourhood. Lemma~\ref{lem:separator} applies with $r(t)=\varepsilon(1+t^2)^{-3}$.
\end{proof}

The proof never uses where the zeros lie inside the open strip; the Euler product enters only through zero-freeness of $\re s\ge1$. For the Davenport--Heilbronn function, which has a functional equation of the same kind but no Euler product, the hypotheses fail for every vertical line, since it has zeros with $\re s>1$ (Saias--Weingartner \cite{SaiasWeingartner}) and hence, by reflection, zeros with $\re s<0$ (\note{24\_} Proposition 24.6); whether the splitting itself fails for it is open.
Two structural facts complete the picture (\lbl{GMS}; \note{24\_} §§2.3--2.4): the derived separation $\operatorname{RHom}(Q,Q_+)\simeq0\simeq\operatorname{RHom}(Q_+,Q)$ of $Q$ and its normal twist $Q_+=\B/I_+$ (jets at $\Zs+1$) is the general vanishing of Ext for a central element acting as $0$ on one module and as $1$ on the other; and $\B$ itself does not split along the two lines, since $Q_+$ is torsion and $\B$ torsion-free, and an entire $c$ with $c^2=c$ is constant.

\subsection{The defect ratio}

For $r>1$ the programme's defect at scale $r$ has eigenvalue $d_r(\rho)=r^{\bar\rho}-r^{1-\rho}=\cl{r^\rho-r^{\rho^\#}}$ at a zero $\rho$: a zero minus its reflected partner.

\begin{lemma}[the exact range of the defect ratio]\label{lem:defect}
Let $r\ne t>1$, $0<\beta<1$, $\beta\ne\tfrac12$, $x=\beta-\tfrac12$. Then $|d_r(\rho)/d_t(\rho)|=\sqrt{r/t}\,\sinh(x\log r)/\sinh(x\log t)$; it is strictly monotone in $|x|$ on $(0,\tfrac12)$, and its values fill exactly the open interval between $\sqrt{r/t}\,\log r/\log t$ (the limit $x\to0$) and $(r-1)/(t-1)$ (the value at $|x|=\tfrac12$).
\end{lemma}
\status{Proved here (claude-ab, \note{24\_} Lemma 24.5; it sharpens the programme's bound \lbl{AST6.2}). Checked on a grid. Refereed (eighth pass). A statement about admissible positions, not about zeros.}

\begin{proof}
$r^\beta-r^{1-\beta}=\sqrt r(r^x-r^{-x})=2\sqrt r\sinh(x\log r)$. The derivative of $\log(\sinh(ax)/\sinh(bx))$ is $a\coth(ax)-b\coth(bx)$, and $c\mapsto c\coth(cx)$ is strictly increasing for $x>0$ (its derivative is $(\tfrac12\sinh2cx-cx)/\sinh^2cx>0$). As $x\to0$ the ratio tends to $\sqrt{r/t}\,a/b$; at $|x|=\tfrac12$, $\sinh(\tfrac12\log r)=(r-1)/(2\sqrt r)$.
\end{proof}

Two further exact statements belong here (register S38, S39; \note{22\_} Propositions 22.1--22.3, 22.5): the Haar and Lebesgue normalisations of ``running the clock backwards'', $g\mapsto\cl{g(1/u)}$ and $g\mapsto\cl{g(1/u)}/u$, have Mellin shadows $s\mapsto-\bar s$ and $s\mapsto1-\bar s$ and differ by the factor $u$; and, paired with their one-line shifts, $\re s=\tfrac12$ and $\re s=-\tfrac12$ are the two sheets of a connected double cover of the equator of $\C P^1$ whose deck transformation is the central $-1$ of $SU(2)$, the author's Hopf fibre and the shift circle being fibres of the left and the right Hopf fibrations.
